\section{Phase Space Factorization and Mapping}

A central ingredient of the Catani--Seymour subtraction method is the
exact factorization of the $(m+1)$-particle phase space into an
$m$-particle phase space and an unresolved emission measure.

\subsection{Phase-space definition}

The $(m+1)$-particle phase space is defined as
\begin{align}
d\Phi_{m+1}(P; p_1,\dots,p_{m+1}) =
\left( \prod_{i=1}^{m+1} \frac{d^3 p_i}{(2\pi)^3 2E_i} \right)
(2\pi)^4 \delta^{(4)}\!\left(P - \sum_{i=1}^{m+1} p_i \right).
\end{align}

For subtraction to work, this phase space must be rewritten in terms of
a reduced $m$-particle configuration.

\subsection{Factorization}

For each dipole $(ij,k)$, the phase space can be factorized as
\begin{align}
d\Phi_{m+1}(p_1,\dots,p_{m+1}) =
d\Phi_m(\tilde{p}_1,\dots,\tilde{p}_m)
\, [dp_i(\tilde{p}_{ij}, \tilde{p}_k)],
\end{align}
where:
\begin{itemize}
    \item $\tilde{p}$ denotes mapped momenta,
    \item $[dp_i]$ is the dipole phase-space measure associated with the
    unresolved emission.
\end{itemize}

This factorization is exact and includes a non-trivial Jacobian.

\subsection{Momentum mapping}

The mapping $(p_i, p_j, p_k) \to (\tilde{p}_{ij}, \tilde{p}_k)$ is constructed such that:
\begin{align}
\tilde{p}_{ij} + \tilde{p}_k = p_i + p_j + p_k,
\end{align}
and:
\begin{align}
\tilde{p}_{ij}^2 = p_{ij}^2, \quad \tilde{p}_k^2 = p_k^2.
\end{align}

A commonly used mapping (final-final case) is:
\begin{align}
\tilde{p}_{ij}^\mu &= p_i^\mu + p_j^\mu - \frac{y_{ij,k}}{1 - y_{ij,k}} p_k^\mu, \\
\tilde{p}_k^\mu &= \frac{1}{1 - y_{ij,k}} p_k^\mu,
\end{align}
with:
\begin{align}
y_{ij,k} = \frac{p_i \cdot p_j}{p_i \cdot p_j + p_j \cdot p_k + p_k \cdot p_i}.
\end{align}

\subsection{Dipole phase-space measure}

The unresolved phase-space factor can be written as
\begin{align}
[dp_i(\tilde{p}_{ij}, \tilde{p}_k)] =
\frac{(1 - y_{ij,k})^{d-3}}{16\pi^2}
\, (p_{ij,k}^2)^{1-\epsilon}
\, dy_{ij,k} \, dz_i \, d\phi,
\end{align}
where:
\begin{itemize}
    \item $y_{ij,k}$ controls the invariant mass of the splitting,
    \item $z_i$ is the energy fraction,
    \item $\phi$ is the azimuthal angle.
\end{itemize}

\subsection{Role in the implementation}

This factorization allows the real-emission contribution to be rewritten as:
\begin{align}
\int d\Phi_{m+1} \, d\sigma^R =
\int d\Phi_m \int [dp_i] \, d\sigma^R.
\end{align}

The subtraction term uses the same mapping, ensuring that:
\begin{itemize}
    \item both $d\sigma^R$ and $d\sigma^A$ are evaluated consistently,
    \item the difference remains finite,
    \item the integrated dipoles can be computed analytically.
\end{itemize}